% --- PAGE 20: HALF-TITLE SYSTEM ANALYSIS ---
\vspace*{\fill}
\begin{center}
  {\fontsize{18pt}{24pt}\selectfont\bfseries SYSTEM ANALYSIS\par}
\end{center}
\vspace*{\fill}
\newpage

% --- PAGE 21: CHAPTER 3 SYSTEM ANALYSIS ---
% --- CHAPTER 3 SYSTEM ANALYSIS ---
\begin{center}
  {\fontsize{16pt}{20pt}\selectfont\bfseries CHAPTER 3\par}
  \vspace{0.3cm}
  {\fontsize{16pt}{20pt}\selectfont\bfseries SYSTEM ANALYSIS\par}
\end{center}
\vspace{0.4cm}

{\fontsize{14pt}{18pt}\selectfont\bfseries 3.1 EXISTING SYSTEM\par}
\vspace{0.3cm}

{\fontsize{12pt}{17.5pt}\selectfont
In prevailing academic operating models, departmental task governance, programming assignment verification, and student career readiness tracking are conducted through disjointed, ad-hoc mechanisms. Faculty members, lab instructors, and student coordinators rely on a patchwork of general-purpose consumer tools, including Google Forms, WhatsApp broadcast channels, shared Google Sheets, and printed observation notebooks. While these tools offer initial convenience, their complete lack of architectural integration introduces systemic operational failures:

\begin{enumerate}[noitemsep,topsep=2pt,leftmargin=*]
  \item \textbf{Deliverable Fragmentation and Zero Auditability:} Student submissions are scattered across personal Google Drives, individual chat messages, and paper files. There is no immutable record of when a student submitted code, who verified it, or the historical feedback provided, making formal NBA and NAAC accreditation audits exceedingly cumbersome.
  \item \textbf{Absence of Native Sandboxed Code Compilation:} Conventional Learning Management Systems (LMS) treat source code files as passive binary attachments. Evaluators must manually download, extract, configure compiler dependencies, and execute code on local workstations, creating multi-week grading backlogs.
  \item \textbf{Algorithmic Telemetry Blind Spot:} Students solve challenges on LeetCode and GitHub independently without faculty having real-time visibility into solve velocity, problem difficulty breakdowns, or commit streaks.
  \item \textbf{Unreliable Alerting and High Deadline Miss Rates:} Deadlines communicated via broadcast emails suffer from low open rates ($< 25\%$), leading to frequent missed submissions and administrative friction.
\end{enumerate}

\vspace{0.4cm}
{\fontsize{14pt}{18pt}\selectfont\bfseries 3.2 EXISTING SYSTEM DISADVANTAGES\par}
\vspace{0.3cm}

A deep architectural audit reveals several critical vulnerabilities in the existing operating model:
\begin{enumerate}[noitemsep,topsep=2pt,leftmargin=*]
  \item \textbf{Zero Automated Disaster Recovery:} Legacy departmental spreadsheets and local databases lack automated, offsite backup mechanisms. A server hardware crash or accidental file deletion results in permanent, irrecoverable loss of student submission histories.
  \item \textbf{Single-Tier Insecure Authentication:} Existing portals employ flat access controls without multi-role segregation, exposing critical grade data to unauthorized manipulation. There is no cryptographically verified one-time password (OTP) mechanism for password recovery.
  \item \textbf{Unoptimized Storage Pipelines:} Submissions are uploaded directly to application servers as uncompressed local files, exhausting server disk capacity and risking unauthorized directory traversal.
  \item \textbf{Absence of Anti-Cheat Proctoring:} Online tests conducted via Google Forms provide zero protection against tab switching, screen sharing, or impersonation during technical screening tests.
\end{enumerate}

\vspace{0.4cm}
{\fontsize{14pt}{18pt}\selectfont\bfseries 3.3 PROPOSED SYSTEM\par}
\vspace{0.3cm}

The proposed \textbf{VSBEC IT Vault} platform overcomes every limitation of the legacy setup by establishing an integrated, cloud-native architecture engineered specifically for the Department of Information Technology at V.S.B. Engineering College.

The platform is structured across five synchronized engineering pillars:
First, a responsive single-page web client built with React 19, Vite 6, and Tailwind CSS provides customized role dashboards for Supreme Admins, HOD, Faculty Class Advisors, Student Coordinators, Students, and Industry Partners.
Second, a resilient Node.js 20 and Express 4.21 application layer enforces strict Role-Based Access Control (RBAC) via cryptographically signed JSON Web Tokens (JWT) and Bcrypt hashing.
Third, an enterprise PostgreSQL database hosted on Supabase utilizes dynamic transaction-mode connection pooling on port 6543, eliminating serverless connection exhaustion.
Fourth, a multi-language sandboxed compiler engine isolates student code executions in C, C++, Java 17, and Python 3 inside secure process jails with 3000ms execution timeout traps and automatic failover to the Judge0 CE cloud API.
Fifth, an automated disaster recovery workflow running on GitHub Actions executes twice daily (07:50 AM IST and 11:55 PM IST), exporting complete student progress reports, database JSON snapshots, and syncing code to dual Git remotes.

\vspace{0.4cm}
{\fontsize{14pt}{18pt}\selectfont\bfseries 3.4 PROPOSED SYSTEM ADVANTAGES AND COMPARISON MATRIX\par}
\vspace{0.3cm}

The proposed system delivers measurable operational improvements across departmental governance, academic velocity, and assessment integrity. Table 3.1 details a comprehensive feature-by-feature comparison between traditional academic tools and the verified capabilities of the VSBEC IT Vault platform.\par}

\vspace{0.3cm}
\begin{table}[H]
  \centering
  \fontsize{9.5pt}{13.5pt}\selectfont
  \begin{tabularx}{\textwidth}{|p{3.2cm}|p{5.8cm}|L|}
    \hline
    \textbf{Dimension / Feature} & \textbf{Traditional Academic Systems (LMS / Forms / Sheets)} & \textbf{Proposed VSBEC IT Vault Architecture} \\
    \hline
    \textbf{Role-Based Access Control (RBAC)} & Flat single-tier or two-tier permissions; no peer coordinator verification layer. & Granular 6-role hierarchy (\texttt{SUPREME\_ADMIN}, \texttt{HOD}, \texttt{ADVISOR}, \texttt{COORDINATOR}, \texttt{STUDENT}, \texttt{INDUSTRY}). \\
    \hline
    \textbf{Disaster Recovery \& Auto Backup} & Manual file downloads; zero automated offsite backup or version-controlled sync. & Automated GitHub Actions cron (07:50 AM \& 11:55 PM IST), JSON DB snapshots, and dual-repo sync. \\
    \hline
    \textbf{Transactional Email \& OTP Engine} & Generic unauthenticated SMTP relays with high spam delivery and no load balancing. & High-availability multi-node Brevo pool (up to 5 nodes) with Resend fallback and 6-digit OTP verification. \\
    \hline
    \textbf{Code Evaluation Runtime} & Passive file storage; faculty must download and compile code manually. & Sandboxed execution jail (C, C++, Java, Python) with 3000ms timeout traps and Judge0 cloud failover. \\
    \hline
    \textbf{Competitive Coding Telemetry} & Zero integration with LeetCode or GitHub; manual student self-reporting. & Nightly GraphQL/REST daemons harvesting solved counts, streaks, and GitHub commit histories. \\
    \hline
    \textbf{Asset \& Media Pipeline} & Uncompressed server disk storage or personal Google Drive links. & Cloudinary HTTPS storage with signed multipart uploads and automatic stale screenshot purging. \\
    \hline
    \textbf{Multi-Channel Real-Time Alerts} & Mass broadcast emails with $< 25\%$ open rates and high deadline miss rates. & Event-driven tri-channel alerts across Telegram bot webhooks, VAPID Web Push, and transactional email. \\
    \hline
    \textbf{Integrity \& Proctoring Engine} & None; open browser tabs permit unrestricted external collaboration. & Client-side proctoring with webcam PIP streams, HTML5 fullscreen enforcement, and tab-switch traps. \\
    \hline
  \end{tabularx}
  \vspace{4pt}
  {\centering\fontsize{10.5pt}{14pt}\selectfont Table 3.1 Comparative Analysis: Traditional Systems vs. VSBEC IT Vault\par}
\end{table}
\newpage
